\section{The full relative trace comparison}
\label{app:formal-trace-comparison}

We prove Proposition~\ref{prop:formal-trace-comparison}, including its
compatibility with the relative integration homotopy.  This is the
comparison used in Section~\ref{sec:algebraicity} to identify the
shifted symplectic structure on the punctured filling with the integrated
trace on the edge complex.

We retain the closed-form complex and total signs of
Section~\ref{subsec:formal-trace-statement}.  There are three
differentials: the source differential \(d_X\), the coefficient
cdga differential \(Q\), and coordinate de Rham differentiation
\(\delta_R\).  On coordinate forms \(Q\) acts by
\(\mathcal L_Q\).  Source and simplex forms contribute to internal
degree but have exterior weight zero; the coordinate forms supply
exterior weight.  Coefficient completion and the product over exterior
weights are both retained throughout the argument.

\begin{proof}[Proof of Proposition~\ref{prop:formal-trace-comparison}]
\noindent\textbf{1. The target cocycle.}\quad
We first construct the target form on a presentation of the formal
classifying stack.  Let \(\widehat G\) be the formal group of \(G\)
at the identity.  Here the pointed completion of \(BG\) means the
deformation functor with a specified identification of its reduction
with the trivial object; its automorphisms reduce to the identity.
This pointed completion is \(B\widehat G\), presented by the group nerve
\(\widehat G^\bullet\).  Its term \(\widehat G^n\) is formally smooth,
with completed coordinate algebra
\[
                  R_n=\mathbb C[[x_1,\ldots,x_{n\dim G}]].
\]
Continuous coordinate forms on these algebras therefore compute
derived forms on the nerve.

We use backward transition elements for the group nerve.  More explicitly,
if \(U_{ij}\) is the forward transport from vertex \(i\) to vertex \(j\)
for the connection \(d+a\), then
\(U_{ik}=U_{jk}U_{ij}\).  The nerve labels
\(h_{ij}=U_{ij}^{-1}\) therefore satisfy
\(h_{ik}=h_{ij}h_{jk}\).  For a left Maurer--Cartan form
\(a=\gamma^{-1}d\gamma\), these labels are
\(h_{ij}=\gamma(i)^{-1}\gamma(j)\), whereas the forward transport is
\(U_{ij}=\gamma(j)^{-1}\gamma(i)\).  Thus Getzler's interval labeled
by \(\exp(x)\) has form \(x\,dt\), and its forward holonomy is
\(\exp(-x)\), as in the edge construction.  This is the opposite-arrow
convention for describing the same local system; inversion is not being
regarded as a homomorphism from a noncommutative group to itself.  All
based holonomy markings below remain the forward transport matrices.

For an ordinary nilpotent ideal \(\mathfrak m\), the Lie algebra
\(\mathfrak g\otimes\mathfrak m\) is nilpotent.  Getzler's simplicial
Maurer--Cartan space is
\[
 \MC_\bullet(\mathfrak g\otimes\mathfrak m)
 =\MC\bigl(\mathfrak g\otimes\mathfrak m
                     \otimes\Omega_{\mathrm{poly}}^*(\Delta^\bullet)\bigr).
\]
Its Dupont-gauge subspace \(\gamma_\bullet\), defined using the
simplicial de Rham contraction of \cite{Dupont1976}, is the nerve of the
Campbell--Hausdorff group in this degree-zero case, and its inclusion
in \(\MC_\bullet\) is a natural homotopy equivalence
\cite[\S1, Theorem~5.4, and Corollary~5.11]{Getzler2009}.
Applying this construction to the universal \(n\)-tuple, first at
each finite augmentation order and then taking the inverse limit,
gives compatible flat simplices
\[
 a_n\in
 \bigl(\mathfrak g\otimes\mathfrak m_{R_n}
          \widehat\otimes\Omega_{\mathrm{poly}}^*(\Delta^n)\bigr)^1.
\]
The compatibility here is that of Getzler's simplicial construction,
including every face and degeneracy.

Differentiate only in the \(R_n\) directions and set
\begin{equation}
\label{eq:nerve-trace-cocycle}
                       \beta_n
                  =\tfrac12\kappa(\delta_{R_n}a_n,\delta_{R_n}a_n).
\end{equation}
Polynomial simplex forms have internal degree but exterior weight
zero in this expression.  Differentiating the Maurer--Cartan equation
gives the covariant-closedness equation for \(\delta_{R_n}a_n\).
Invariance of \(\kappa\) cancels the two bracket terms and yields
\[
                         d_\Delta\beta_n=0,
                         \qquad \delta_{R_n}\beta_n=0.
\]
These identities and simplicial compatibility make \(\beta_\bullet\)
a Thom--Whitney descent cocycle of internal degree two and exterior
weight two on the smooth formal nerve.  All its higher closed-form
components are zero.  Equivalently, integration over the simplices
gives a normalized descent cocycle; its descent differential is zero
by Stokes' formula.

\smallskip\noindent\textbf{2. The canonical closed-form class.}\quad
We identify its class with the restriction of \(\omega_\kappa\).
The cotangent complex of \(B\widehat G\) is the coadjoint bundle
\(\mathfrak g^*[-1]\).  Consequently
\[
 R\Gamma(B\widehat G,\mathop{\bigwedge}\nolimits^p
                  \mathbb L_{B\widehat G})\in D^{\geq p}.
\]
Indeed, \(\bigwedge^p\mathbb L\) is the bundle
\(\Sym^p(\mathfrak g^*)[-p]\), and derived invariants of a bundle in
degree zero have no negative cohomology.  Applied to the exterior
weight filtration of the closed two-form complex, this connectivity
gives
\begin{equation}
\label{eq:formal-bg-form-uniqueness}
 H^2\operatorname{Cl}^2(B\widehat G)
       \simeq \Sym^2(\mathfrak g^*)^{\widehat G},
 \qquad H^1\operatorname{Cl}^2(B\widehat G)=0.
\end{equation}
For the \(j\)-th higher component, the relevant internal degree would
be \(2-j\), whereas weight \(2+j\) starts in degree \(2+j\).
This is a statement about the derived form complexes and their
classes, independent of the terms in a chosen descent resolution.

Linearizing the nerve at the identity gives the Dold--Kan nerve of
\(\mathfrak g[1]\).  Formula~\eqref{eq:nerve-trace-cocycle}
linearizes to its quadratic two-form with associated symmetric
bilinear form \(\kappa\).  Here we use the usual normalized simplex
integration and shuffle convention: the two shuffles in the
polarization cancel the factor \(1/2\).
This is precisely the underlying bilinear form of
\(\omega_\kappa\) \cite[\S1.2, classifying stacks]{PTVV2013}.
Equation~\eqref{eq:formal-bg-form-uniqueness} identifies their full
closed-form classes, and makes the comparison homotopy unique up
to homotopy.

\smallskip\noindent\textbf{3. Derived coefficients.}\quad
We now give the comparison on derived coefficients, including its map on
linear fibers.  By an Artin test algebra here we mean a nonpositive local
augmented cdga \(B\) over \(\mathbb C\), with nilpotent augmentation
ideal \(\mathfrak m_B\).  Set
\[
 L_B=\mathfrak g\otimes\mathfrak m_B,
 \qquad \Omega_q=\Omega^*_{\mathrm{poly}}(\Delta^q),
 \qquad H_q(B)=\exp Z^0(L_B\otimes\Omega_q).
\]
The exponential uses the finite Campbell--Hausdorff series.  These groups
form a simplicial group \(H_\bullet(B)\).  They compute the derived
evaluation of \(\widehat G\) on \(B\): in exponential coordinates a
map from its completed smooth coordinate algebra sends each coordinate
to a degree-zero cycle in \(\mathfrak m_B\otimes\Omega_q\).  Equivalently
one uses the nonpositive truncation of this polynomial simplicial frame.
Consequently the diagonal of the bisimplicial nerve
\(N_nH_q(B)\) computes \(B\widehat G(B)\).

Apply the ordinary degree-zero Getzler construction to an \(n\)-tuple
in \(H_q(B)\).  The inclusion of the degree-zero Lie algebra
\(Z^0(L_B\otimes\Omega_q)\) into \(L_B\otimes\Omega_q\) gives a
flat element
\[
 A_{n,q}\in
 \MC\bigl(L_B\otimes\Omega_q\otimes\Omega_n\bigr).
\]
The coefficient cycles have zero total differential, so this is the
Maurer--Cartan equation for the full displayed tensor product.
Getzler's compatibility in \(n\) and functoriality in the coefficient
Lie algebra give compatibility in both simplicial directions, including
all faces and degeneracies.  Restriction to the diagonal simplex gives
a natural simplicial map
\begin{equation}
\label{eq:explicit-dg-nerve-comparison}
 \Phi_B:\operatorname{diag}N_\bullet H_\bullet(B)
                  \longrightarrow\MC_\bullet(L_B),
 \qquad
 (\Phi_B)_q=\Delta_q^*A_{q,q}.
\end{equation}

Here is a direct verification that this particular map is an equivalence.
Consider a central step \(B\twoheadrightarrow B'\) in the filtration
by powers of \(\mathfrak m_B\), with kernel \(I\) satisfying
\(\mathfrak m_BI=0\), and write \(V=\mathfrak g\otimes I\).
The map \(Z^0(L_B\otimes\Omega_q)\to
Z^0(L_{B'}\otimes\Omega_q)\) is surjective.  Indeed, lift a closed
element to \(w\).  Its differential \(Dw\) is a closed degree-one
element of \(V\otimes\Omega_q\); the polynomial Poincar\'e lemma and
nonpositivity give
\(H^1(V\otimes\Omega_q)=H^1(V)=0\).  Subtracting a degree-zero
primitive of \(Dw\) gives a closed lift.  Thus the induced map of
simplicial groups is degreewise surjective, with additive kernel
\[
                         K_q=Z^0(V\otimes\Omega_q).
\]
Its classifying-space homotopy fiber is \(BK_\bullet\).
On the Maurer--Cartan side the surjection is a Kan fibration with fiber
\(\MC_\bullet(V)\), by \cite[Proposition~4.7]{Getzler2009}.
We use homotopy fibers of the classifying spaces here; no Kan-fibration
assertion about an arbitrary diagonal map is needed.

The restriction of \(\Phi_B\) to these fibers can be checked explicitly.
For the abelian complex \(V\), put
\[
 C_q=(V\otimes\Omega_q)^0,
 \qquad M_q=Z^1(V\otimes\Omega_q)=\MC_q(V).
\]
There is a degreewise exact sequence of simplicial vector spaces
\[
                    0\longrightarrow K_\bullet
                    \longrightarrow C_\bullet
                    \xrightarrow{D}M_\bullet\longrightarrow0.
\]
Let \((EK)_q=K_q^{q+1}\) be the homogeneous bar construction, with
diagonal translation action of \(K_q\).  Its quotient is the diagonal
inhomogeneous nerve \(BK\).  The map
\[
 \Psi_q:(EK)_q\longrightarrow C_q,
 \qquad
       (k_0,\ldots,k_q)\longmapsto\sum_{j=0}^q t_j k_j
\]
is simplicial and restricts to the identity on the diagonal copy of
\(K_q\).  Since \(Dk_j=0\), its quotient map is
\(D\Psi_q=\sum_jdt_j k_j\).  In the inhomogeneous coordinates
\(x_i=k_i-k_{i-1}\), this is exactly the abelian Getzler simplex
after the diagonal restriction in
\eqref{eq:explicit-dg-nerve-comparison}.  Thus we have a commuting
diagram of exact simplicial vector spaces
\[
\begin{array}{ccccccccc}
0&\longrightarrow&K&\longrightarrow&EK&\longrightarrow&BK&\longrightarrow&0\\
 &&\Vert&&\downarrow\Psi&&\downarrow\Phi_V\\
0&\longrightarrow&K&\longrightarrow&C&\xrightarrow{D}&M&\longrightarrow&0.
\end{array}
\]
Both middle terms are contractible.  This can be seen without choosing
a contraction of \(V\): for a simplicial interval parameter
\(\sigma:[q]\to[1]\), multiplication by
\(\sum_{\sigma(j)=1}t_j\) gives a simplicial homotopy from zero to
the identity on \(C_q\).  On \((EK)_q\), use
\((k_j)_j\mapsto(\mathbf1_{\sigma(j)=1}k_j)_j\).
The connecting maps of the two exact sequences therefore identify
\(\Phi_V\) with the identity suspension of \(K\).  In particular it
is a weak equivalence, and under the compatible integration and
suspension identifications it is the canonical map from the bar model
of \(V\) to the model of \(V[1]\).  This statement uses those
connecting-map conventions, not an unsigned equality of raw simplex
integrals.  Induction over the finite central filtration, starting at
\(B=\mathbb C\), now proves that \(\Phi_B\) is an equivalence.
The construction is natural in \(B\); the usual nilpotent
Maurer--Cartan invariance then also gives its compatibility with
quasi-isomorphisms.  This is a proof of the dg extension, rather than
an assertion that the ordinary coefficient-ring statement of
\cite[Proposition~4.9]{Getzler2009} already states that extension.

\smallskip\noindent\textbf{4. Compatibility of the closed cocycle.}\quad
We finally check the closed cocycle under this very comparison.  Do this
before simplex integration, in the Thom--Whitney mixed complexes, and
use a completed cofibrant coefficient resolution when computing derived
forms.  A tuple in \(H_q(B)\) gives the substitution of the universal
nerve coordinates by closed elements of
\(\mathfrak m_B\otimes\Omega_q\).  Pull back
\(\beta_n=\tfrac12\kappa(\delta a_n,\delta a_n)\) by that
substitution, and apply the coefficient morphism
\[
 \operatorname{DR}(B\otimes\Omega_q)
             \longrightarrow\operatorname{DR}(B)\otimes\Omega_q.
\]
As throughout this proof, this notation means derived de Rham objects
computed on the chosen resolutions.  The morphism puts the polynomial
simplex forms in exterior weight zero.  In particular coordinate
differentiation still differentiates every \(B\)-dependent coefficient
of the logarithms of the tuple; those coefficients are not held
constant.  Naturality of coordinate differentiation gives precisely
\[
                \frac12\kappa(\delta_B A_{n,q},
                                      \delta_B A_{n,q}).
\]
Both simplex directions have exterior weight zero.  Diagonal restriction
commutes with their internal differential and with \(\delta_B\), with
the total signs already specified.  Consequently
\[
 \Delta_q^*\!\left(\frac12\kappa(\delta_B A_{q,q},
                                         \delta_B A_{q,q})\right)
       =\frac12\kappa\bigl(\delta_B(\Phi_B)_q,
                          \delta_B(\Phi_B)_q\bigr)
\]
as full mixed cocycles.  The higher closed-form components of these
representatives are all zero; no exterior weights have been dropped.
Faces, degeneracies, coefficient maps, and their homotopies commute with
this equality.  The standard Thom--Whitney-to-normalized comparison
therefore gives equality of the full closed-form classes.  We do not
identify integration over a product of simplices with integration over
its diagonal.  This establishes the derived pullback of the class
already constructed on the smooth nerve, together with its natural
comparison homotopy.  Passing continuously through the nilpotent
coefficient quotients gives the same assertion for a completed semifree
\(R\), where continuous coordinate forms compute the derived forms.

\smallskip\noindent\textbf{5. Source evaluation and relative integration.}\quad
We next carry out evaluation on the source.  Choose a finite
triangulation \(K\) of \(X\).  A compatible family of flat simplices
\[
 a_\sigma\in
 \MC\bigl(\mathfrak g\otimes\mathfrak m_R
                    \widehat\otimes\Omega_{\mathrm{poly}}^*(\Delta^{\dim\sigma})
       \bigr),\qquad \sigma\in K,
\]
is a simplicial map from \(K\) to the Maurer--Cartan model of
\(B\widehat G\).  Including simplicial test parameters gives its
families and all their homotopies.  Evaluation is the evaluation
of these simplicial maps.  Pulling back the target cocycle therefore
assigns to each \(\sigma\) the expression
\(\tfrac12\kappa(\delta_Ra_\sigma,\delta_Ra_\sigma)\).

Before totalizing this diagram, apply the coefficient map
\[
 \kappa_{\Spf R,X_B}:
 \operatorname{DR}(\Spf R\times X_B)
       \longrightarrow
 \operatorname{DR}(\Spf R)\widehat\otimes
                         R\Gamma(X_B,\mathcal O_{X_B})
\]
used in \cite[\S2.1, before Definition~2.3]{PTVV2013}.
It keeps source functions and cochains in exterior weight zero.
Thus the polynomial simplex forms above belong to the source
cochain factor, while \(\delta_R\) supplies every exterior weight
in the coefficient form.  Totalizing afterwards produces
exactly \(\Theta_X\) in
\eqref{eq:formal-trace-representative}.  We have proved the equality
of full closed-form classes
\[
     \kappa_{\Spf R,X_B}
                  \operatorname{ev}_{a_X}^*[\omega_\kappa]
                          =[\Theta_X].
\]
All these operations are performed weight by weight and are
continuous in the coefficient ideal.  They then define the
displayed completed closed-form complex.

For a smooth triangulated manifold, the same statement uses smooth
forms through the multiplicative comparison
\[
 \Omega_{\mathrm{sm}}^*(X)
       \longrightarrow\Omega_{\mathrm{ps}}^*(K)
       \longleftarrow\Omega_{\mathrm{poly}}^*(K),
\]
where the middle term consists of compatible piecewise smooth forms.
Both arrows are quasi-isomorphisms and commute with restriction.
Integration of forms over the oriented simplices is the same on
all three models, including their boundaries.  The construction
therefore commutes with the orientation cochain maps, as well as
with any multiplicative replacement carrying those maps.

Finally let \(\partial Z=S\).  Write \(I_Z\) and \(I_S\) for
integration over compatible fundamental chains.  With
\(\partial[Z]=[S]\) and the suspension convention for transgression,
the equations in~\eqref{eq:formal-trace-representative} give
\[
 h_Z=I_Z\Theta_Z,\qquad
 \mathcal L_Qh_Z=I_S\Theta_S,\qquad \delta_Rh_Z=0.
\]
Thus \(h_Z\), of internal degree \(-1\), is the relative
closed-form homotopy from zero to the pulled-back boundary form.
This is the homotopy obtained by applying the orientation
construction to evaluation
\cite[Definition~2.6 and Claim~2.7]{Calaque2015}.
The comparison homotopy on \(B\widehat G\) is carried through the
same evaluation and relative-chain maps.  It consequently compares
these relative homotopies together with their endpoints.

On the sphere condition \(a_S=\nu b\), every term of
\(\kappa(\delta_Ra_S,\delta_Ra_S)\) contains \(\nu\wedge\nu=0\).
Its full closed-form representative and its chosen path are
therefore zero.  Subtracting this zero path in the shifted
Lagrangian intersection construction leaves \(h_Z\), now closed
of degree \(-1\), as asserted.
\end{proof}
